\documentclass[a4paper,11pt]{article}
\usepackage[utf8]{inputenc}
\usepackage{geometry}
\usepackage{xcolor}
\usepackage{listings}
\usepackage{hyperref}
\usepackage{graphicx}
\usepackage{amsmath}
\usepackage{amssymb}
\usepackage{tcolorbox}
\usepackage{float}
\usepackage{booktabs}
\usepackage{tikz}
\usetikzlibrary{arrows.meta, positioning, shapes.geometric, matrix}

% ── Page geometry ──────────────────────────────────────────────────────────────
\geometry{top=2.5cm, bottom=2.5cm, left=2.5cm, right=2.5cm}

% ── Custom colours ─────────────────────────────────────────────────────────────
\definecolor{ibmblue}{RGB}{0,45,156}
\definecolor{ibmred}{RGB}{250,75,76}
\definecolor{codegreen}{rgb}{0,0.6,0}
\definecolor{codegray}{rgb}{0.5,0.5,0.5}
\definecolor{codepurple}{rgb}{0.58,0,0.82}
\definecolor{backcolour}{rgb}{0.95,0.95,0.92}
\definecolor{tealcolor}{RGB}{0,128,128}
\definecolor{orangecolor}{RGB}{230,100,0}
\definecolor{goldcolor}{RGB}{180,140,0}

% ── Colour helpers ─────────────────────────────────────────────────────────────
\newcommand{\crd}[1]{\textcolor{ibmred}{#1}}
\newcommand{\cpur}[1]{\textcolor{codepurple}{#1}}
\newcommand{\cblu}[1]{\textcolor{ibmblue}{#1}}
\newcommand{\cgr}[1]{\textcolor{codegreen}{#1}}
\newcommand{\corg}[1]{\textcolor{orangecolor}{#1}}
\newcommand{\cteal}[1]{\textcolor{tealcolor}{#1}}
\newcommand{\cgold}[1]{\textcolor{goldcolor}{#1}}

% ── Listings ───────────────────────────────────────────────────────────────────
\lstdefinestyle{mystyle}{
    backgroundcolor=\color{backcolour},
    commentstyle=\color{codegreen},
    keywordstyle=\color{magenta},
    numberstyle=\tiny\color{codegray},
    stringstyle=\color{codepurple},
    basicstyle=\ttfamily\footnotesize,
    breakatwhitespace=false, breaklines=true,
    captionpos=b, keepspaces=true,
    numbers=left, numbersep=5pt,
    showspaces=false, showstringspaces=false,
    showtabs=false, tabsize=2
}
\lstset{style=mystyle}

% ── tcolorbox environments ─────────────────────────────────────────────────────
\newtcolorbox{studentbox}{colback=blue!5!white,colframe=blue!75!black,
    title=Student Activity, fonttitle=\bfseries}
\newtcolorbox{warningbox}[1][Key Concept]{%
    colback=yellow!10!white,colframe=orange!80!black,
    title={#1}, fonttitle=\bfseries}
\newtcolorbox{mathbox}[1][Derivation]{%
    colback=red!5!white,colframe=red!75!black,
    title={#1}, fonttitle=\bfseries}
\newtcolorbox{defbox}[1][Definition]{%
    colback=teal!5!white,colframe=teal!70!black,
    title={#1}, fonttitle=\bfseries}
\newtcolorbox{finbox}[1][Financial Interpretation]{%
    colback=green!5!white,colframe=green!65!black,
    title={#1}, fonttitle=\bfseries}
\newtcolorbox{bridgebox}{colback=purple!5!white,colframe=purple!70!black,
    title={Optional: connection to Lab 4}, fonttitle=\bfseries}

% ── Math shorthands ────────────────────────────────────────────────────────────
\newcommand{\ket}[1]{|#1\rangle}
\newcommand{\bra}[1]{\langle#1|}
\newcommand{\expect}[1]{\langle #1 \rangle}
\newcommand{\Mod}[1]{\ (\mathrm{mod}\ #1)}
\newcommand{\tr}{\mathrm{Tr}}

\title{\textbf{Laboratory Guide 4 (Finance): Quantum Portfolio Optimization}\\
\large Finance \& Industry Applications of Variational Quantum Algorithms}
\author{Course: Quantum Computing Foundation}
\date{}

\begin{document}

\maketitle
\tableofcontents
\newpage

% ══════════════════════════════════════════════════════════════════════════════
\section{Introduction and Objectives}
% ══════════════════════════════════════════════════════════════════════════════

This laboratory applies a quantum optimization algorithm to a real-world
industrial problem: \textbf{portfolio optimization in quantitative finance}.
Everything it needs beyond Labs 1--2 is explained in this guide. If you did
Lab 4, you have met the same algorithm on a small graph problem (MaxCut), and
the optional boxes point out the connections.

The investment choice is written as a Quadratic Unconstrained Binary
Optimization (QUBO) problem, rewritten as an operator on qubits (an Ising
Hamiltonian), and searched with a short quantum circuit, the Quantum
Approximate Optimization Algorithm (QAOA). The novelty is threefold: (i) the QUBO is derived from real financial data,
(ii) interpreting the output requires both a technical and a business lens,
and (iii) the problem has a natural classical solution that allows direct
benchmarking.

\vspace{0.3cm}
\noindent\textbf{Prerequisites:}
\begin{itemize}
    \item Labs 1--2 (Lab 4 helpful, not required)
    \item Basic probability and statistics: mean, variance, covariance
    \item No prior knowledge of quantitative finance is assumed
\end{itemize}

\noindent\textbf{Learning Objectives:}
\begin{enumerate}
    \item Model a combinatorial portfolio selection problem as a QUBO.
    \item Map the QUBO to a quantum Hamiltonian and build the QAOA circuit.
    \item Run QAOA on a 4-asset universe (AAPL, MSFT, JPM, XOM) and
          verify the result against brute-force.
    \item Interpret results from both a business and a technical perspective.
    \item Understand how the circuit and the classical search space grow
          with $n$, and what is (not yet) known about quantum advantage.
\end{enumerate}

\begin{warningbox}[Why Quantum for Portfolio Optimization?]
For $n = 4$, $K = 2$ there are only 6 portfolios, and a laptop checks
them all instantly. For $n = 50$, $K = 10$ there are
$\binom{50}{10} \approx 10^{10}$: checking every one becomes impractical,
and real desks use classical heuristics instead. QAOA's circuit grows only
as $\mathcal{O}(n^2)$ gates per layer, which is why it is studied for this
problem. Whether it can beat the best classical heuristics is an open
research question. This lab shows what it does and does not deliver today.
\end{warningbox}

% ══════════════════════════════════════════════════════════════════════════════
\section{Part A: Classical Portfolio Theory}
% ══════════════════════════════════════════════════════════════════════════════

\subsection{Mean-Variance Optimization (Markowitz, 1952)}

Let $x \in \mathbb{R}^n$ be a portfolio weight vector with
$\sum_i x_i = 1$ (continuous). The mean-variance problem is:

\begin{equation}
    \min_x\; \lambda\, x^T \Sigma\, x - \boldsymbol\mu^T x,
    \label{eq:mv-cont}
\end{equation}

where $\Sigma \in \mathbb{R}^{n\times n}$ is the asset return
\emph{covariance matrix}, $\boldsymbol\mu \in \mathbb{R}^n$ is the
vector of \emph{expected returns}, and $\lambda > 0$ is the
\emph{risk-aversion} parameter. The first term penalises risk (portfolio
variance), the second rewards expected return.

\begin{defbox}[Risk-Aversion Parameter $\lambda$]
\begin{itemize}
    \item $\lambda \to 0$: maximise returns regardless of risk
          (aggressive investor).
    \item $\lambda \to \infty$: minimise variance regardless of returns
          (defensive investor, cash-like).
    \item In this lab we explore $\lambda \in [0.5,\, 5]$.
\end{itemize}
For this lab we use $\lambda = 1$ as the baseline.
\end{defbox}

\subsection{The Combinatorial Version and its Hardness}

In practice, portfolio managers must select \emph{exactly $K$ assets}
from a universe of $n$ candidates.

\noindent\textbf{Declared simplification:} a portfolio here is a
\emph{selection} --- each chosen asset gets the same amount of money;
choosing weights continuously (the full Markowitz problem) is not modelled.

\noindent The binary version of equation~\eqref{eq:mv-cont} is:

\begin{equation}
    \min_{x \in \{0,1\}^n}\; \lambda\, x^T \Sigma\, x
    - \boldsymbol\mu^T x,
    \quad\text{subject to}\quad \mathbf{1}^T x = K.
    \label{eq:mv-binary}
\end{equation}

For $n=4$, $K=2$ there are only $\binom{4}{2} = 6$ candidates — trivially
enumerable. For $n=50$, $K=10$ there are $\binom{50}{10} \approx 10^{10}$
candidates. For $n=500$, $K=50$: $\approx 2.3 \times 10^{69}$ — far beyond brute-force.
The continuous relaxation (equation~\eqref{eq:mv-cont}) gives the
\emph{efficient frontier} but loses the discreteness that real portfolios
require.

\noindent\textbf{Efficient frontier.} A portfolio is \emph{efficient} if no
other portfolio offers a higher expected return for the same or lower risk.
The efficient ones, drawn in the risk--return plane, form the
\emph{efficient frontier}; every other point is \emph{dominated}. Among our
six 2-asset portfolios (equal weights), four are efficient (MSFT + JPM,
AAPL + JPM, AAPL + XOM, AAPL + MSFT); MSFT + XOM and JPM + XOM are
dominated. Moving $\lambda$ walks the optimum along this frontier.

\subsection{Asset Universe: AAPL, MSFT, JPM, XOM (2020--2024)}

We work with the four assets in Table~\ref{tab:assets}. The covariance
matrix $\Sigma$ and expected returns $\boldsymbol\mu$ are computed from
daily closing prices over 2020--2024 (annualised).

\begin{table}[H]
\centering
\caption{Asset data used in this laboratory (annualised, 2020--2024).}
\label{tab:assets}
\begin{tabular}{clccc}
    \toprule
    Index & Ticker & Sector & $\mu_i$ & $\sigma_i = \sqrt{\Sigma_{ii}}$ \\
    \midrule
    0 & AAPL  & Technology    & 0.340 & 0.290 \\
    1 & MSFT  & Technology    & 0.270 & 0.270 \\
    2 & JPM   & Financials    & 0.150 & 0.240 \\
    3 & XOM   & Energy        & 0.220 & 0.350 \\
    \bottomrule
\end{tabular}
\end{table}

\noindent The annualised covariance matrix (approximate):
\begin{equation}
    \Sigma = \begin{pmatrix}
        \cgr{0.0841} & 0.0588 & 0.0313 & 0.0203 \\
        0.0588 & \cgr{0.0729} & 0.0272 & 0.0170 \\
        0.0313 & 0.0272 & \cgr{0.0576} & 0.0294 \\
        0.0203 & 0.0170 & 0.0294 & \cgr{0.1225}
    \end{pmatrix}.
    \label{eq:sigma}
\end{equation}
The diagonal entries $\Sigma_{ii} = \sigma_i^2$ are the individual asset
variances; the off-diagonal entries capture return correlations.
AAPL--MSFT have the highest co-movement (both tech-sector), while
XOM (energy) is weakly correlated with the others.

% ══════════════════════════════════════════════════════════════════════════════
\section{Part B: QUBO Formulation}
% ══════════════════════════════════════════════════════════════════════════════

\subsection{From Constrained to Unconstrained}

We convert the equality constraint $\sum_i x_i = K$ into a quadratic
penalty added to the objective:

\begin{equation}
    \min_{x\in\{0,1\}^n}\;
    \underbrace{\lambda\, x^T\Sigma x - \boldsymbol\mu^T x}_{\text{cost}}
    \;+\;
    \underbrace{\lambda_\text{pen}\!\left(\sum_{i=1}^n x_i - K\right)^{\!2}}_{\text{budget penalty}},
    \label{eq:qubo-full}
\end{equation}

where $\lambda_\text{pen}$ must be large enough that any infeasible
solution (wrong number of assets) is always worse than any feasible one.
A safe choice is $\lambda_\text{pen} > \max_{x,\text{feasible}} |C(x)|$.

Expanding and collecting, equation~\eqref{eq:qubo-full} takes the
standard QUBO form.

\begin{warningbox}[QUBO Derivation — step by step]
Since $x_i \in \{0,1\}$, we have $x_i^2 = x_i$.
Expand the budget penalty (dropping the constant $\lambda_\text{pen}K^2$):
\[
    \lambda_\text{pen}\!\left(\sum_i x_i - K\right)^{\!2}
    \;\longrightarrow\;
    \lambda_\text{pen}(1-2K)\!\sum_i x_i
    \;+\; 2\lambda_\text{pen}\!\sum_{i<j} x_i x_j.
\]
Adding the financial objective and grouping:
\begin{itemize}
  \item \textbf{Diagonal} ($i=j$): coefficient of $x_i$ is
        $\lambda\Sigma_{ii} - \mu_i + \lambda_\text{pen}(1-2K)$.
  \item \textbf{Off-diagonal} ($i<j$, one occurrence): coefficient of $x_ix_j$ is
        $2\lambda\Sigma_{ij} + 2\lambda_\text{pen}$.
\end{itemize}
Convention note: because the full product $x^TQx = \sum_{i,j}Q_{ij}x_ix_j$ counts
each unordered pair twice, the symmetric matrix entry $Q_{ij}$ ($i\neq j$) in
equation~\eqref{eq:Q} carries \emph{half} of the $i<j$ coefficient above — the
risk part and the penalty part alike — giving
$Q_{ij}=\lambda\Sigma_{ij}+\lambda_\text{pen}$.
With this scaling,
\[
    x^TQx \;=\; C(x) \;+\; \lambda_\text{pen}\!\left(\sum_i x_i - K\right)^{\!2}
    \;-\; \lambda_\text{pen}K^2 .
\]
On a feasible portfolio the penalty vanishes, so $x^TQx = C(x) - \lambda_\text{pen}K^2$:
a constant shift, identical for every feasible $x$, which does not affect the ranking.
Every infeasible portfolio pays a strictly positive penalty — that is what enforces
the budget constraint.
\end{warningbox}

\begin{equation}
    \min_{x\in\{0,1\}^n}\; x^T Q\, x,
    \label{eq:qubo}
\end{equation}

\noindent with:
\begin{equation}
    Q_{ij} = \begin{cases}
        \lambda\,\Sigma_{ii} - \mu_i + \lambda_\text{pen}(1-2K)
        & i = j \\[4pt]
        \lambda\,\Sigma_{ij} + \lambda_\text{pen}
        & i \neq j.
    \end{cases}
    \label{eq:Q}
\end{equation}

\begin{mathbox}[QUBO Matrix for $n=4$, $K=2$, $\lambda=1$, $\lambda_\text{pen}=3$]
Substituting equation~\eqref{eq:sigma} and $\mu = (0.34, 0.27, 0.15, 0.22)$
into equation~\eqref{eq:Q} with $\lambda_\text{pen}(1-2K) = 3(1-4) = -9$:
\[
    Q = \lambda\Sigma - \mathrm{diag}(\boldsymbol\mu)
    + \lambda_\text{pen}
    \begin{pmatrix}
        1-2K & 1 & 1 & 1 \\
        1 & 1-2K & 1 & 1 \\
        1 & 1 & 1-2K & 1 \\
        1 & 1 & 1 & 1-2K
    \end{pmatrix}
\]
\[
    Q = \begin{pmatrix}
        -\crd{9.2559} & 3.0588 & 3.0313 & 3.0203 \\
        3.0588 & -\crd{9.1971} & 3.0272 & 3.0170 \\
        3.0313 & 3.0272 & -\crd{9.0924} & 3.0294 \\
        3.0203 & 3.0170 & 3.0294 & -\crd{9.0975}
    \end{pmatrix}.
\]
The positive off-diagonal terms ($\approx\lambda_\text{pen}$) reflect the budget
penalty, which charges every extra pair of selected assets;
the large negative diagonal terms ($\approx\lambda_\text{pen}(1-2K) = -9$) come from the
linear part of the same penalty, and the small remainder $\lambda\Sigma_{ii} - \mu_i$ on
the diagonal is the finance (own risk minus return).
\end{mathbox}

\subsection{Classical Brute-Force Baseline}

For $n=4$, $K=2$ we can enumerate all $\binom{4}{2}=6$ portfolios
exactly. Table~\ref{tab:bf} shows the financial cost
$C(x)=\lambda x^T\Sigma x - \boldsymbol\mu^T x$ for each feasible
portfolio. (All feasible portfolios satisfy the budget constraint, so the
penalty term is zero and $C(x)$ equals the QUBO cost up to the constant
$\lambda_\text{pen}K^2=12$; only differences matter for ranking.)

\begin{table}[H]
\centering
\caption{Brute-force enumeration of all 2-asset portfolios ($\lambda=1$).
Portfolio vector $x=[x_0,x_1,x_2,x_3]$ shown in AAPL/MSFT/JPM/XOM order.}
\label{tab:bf}
\begin{tabular}{clcc}
    \toprule
    Portfolio vector $x$ & Assets selected & Financial cost $C(x)$ & Rank \\
    \midrule
    $\cblu{1100}$ & \cblu{AAPL + MSFT} & $\cblu{-0.335}$ & \cblu{1 (optimal)} \\
    $1001$ & AAPL + XOM        & $-0.313$ & 2 \\
    $1010$ & AAPL + JPM        & $-0.286$ & 3 \\
    $0101$ & MSFT + XOM        & $-0.261$ & 4 \\
    $0110$ & MSFT + JPM        & $-0.235$ & 5 \\
    $0011$ & JPM  + XOM        & $-0.131$ & 6 \\
    \bottomrule
\end{tabular}
\end{table}

\begin{finbox}
\textbf{Why AAPL + MSFT?} Despite having the highest inter-asset
covariance (both technology sector), AAPL and MSFT also have the
highest expected returns ($0.34$ and $0.27$ respectively). For
$\lambda=1$ the return advantage outweighs the diversification cost.
As $\lambda$ increases (more risk-averse investor), the optimal
portfolio shifts to AAPL + JPM (around $\lambda=2$) and eventually to
MSFT + JPM (by $\lambda=5$). Both shifts bring in JPM, which has the
smallest individual variance ($\Sigma_{\text{JPM,JPM}}=0.0576$).
Note that XOM never enters the optimal portfolio, despite having the
smallest covariances with the other assets
($\Sigma_{\text{MSFT,XOM}}=0.0170$ is the smallest off-diagonal entry
of all): portfolio variance $x^T\Sigma x$ includes the diagonal, and
XOM's own variance ($0.1225$, i.e.\ $\sigma=35\%$) is by far the
largest. Low correlation cannot rescue an asset that is individually
too volatile.
\end{finbox}

% ══════════════════════════════════════════════════════════════════════════════
\section{Part C: Quantum Hamiltonian and QAOA Circuit}
% ══════════════════════════════════════════════════════════════════════════════

\subsection{QUBO $\to$ Ising Hamiltonian}

Substituting $x_i = (1 - Z_i)/2$ into equation~\eqref{eq:qubo}.
Note that the eigenvalue of $Z_i$ on state $|q_i\rangle$ is $+1$ if $q_i=0$ and
$-1$ if $q_i=1$, so $x_i=(1-Z_i)/2$ maps $x_i=0\leftrightarrow Z_i=+1$ and
$x_i=1\leftrightarrow Z_i=-1$ consistently.

\begin{align}
    x^T Q\, x &= \sum_{i,j} Q_{ij} x_i x_j
    = \sum_{i,j} Q_{ij} \frac{(1-Z_i)(1-Z_j)}{4} \notag\\
    &= \frac{1}{4}\sum_{i,j} Q_{ij}
    \bigl(I - Z_i - Z_j + Z_iZ_j\bigr) \notag\\
    &\equiv H_\text{portfolio} \;+\; \text{const.}
    \label{eq:ising-map}
\end{align}

To collect terms, use: (a) $Z_i^2 = I$ (constant, absorbed into ``const'');
(b) the $-Z_i$ and $-Z_j$ pieces with symmetric $Q$ give $-\frac{1}{2}\sum_j Q_{ij}$
as coefficient of $Z_i$; (c) cross terms $Z_iZ_j$ with $i\neq j$ give coefficient
$Q_{ij}/2$ after accounting for the $1/4$ prefactor and the factor 2 from
$\sum_{i\neq j} = 2\sum_{i<j}$ (since $Q$ is symmetric).
Collecting terms, the cost Hamiltonian is:

\begin{equation}
    H_\text{portfolio}
    = \sum_{i<j} \frac{Q_{ij}}{2}\, Z_iZ_j
    + \sum_i \left[-\frac{1}{2}\sum_j Q_{ij}\right] Z_i
    + \text{const.}
    \label{eq:H-portfolio}
\end{equation}

A sum of $Z$ and $ZZ$ terms that plays the role of an energy is called an
\textbf{Ising Hamiltonian} (after a physics model of magnets). The cheapest
portfolio is its lowest-energy basis state, the \textbf{ground state}:
minimising cost and finding the ground state are the same task. Because the
substitution drops a constant, energies of $H_\text{portfolio}$ are the QUBO
costs shifted by a fixed amount (here $9.2294$): the best cost $-12.3354$
corresponds to the ground-state energy $-3.1060$.

\begin{bridgebox}
The structure of $H_\text{portfolio}$ is a sum of $Z_iZ_j$ and $Z_i$
Pauli operators. MaxCut in Lab 4 has only $Z_iZ_j$ terms, all with the
same coefficient; the circuit is built from the same blocks:
\begin{itemize}
    \item $e^{-i\gamma J_{ij} Z_i Z_j}$: CNOT\,--\,$R_z(2\gamma J_{ij})$\,--\,CNOT
    \item $e^{-i\gamma h_i Z_i}$: $R_z(2\gamma h_i)$ directly
    \item $e^{-i\beta H_B}$: $R_x(2\beta)$ on all qubits
\end{itemize}
Here the coefficients $J_{ij} = Q_{ij}/2$ and $h_i$ come from the financial
data, and the single-$Z$ terms are new.
\end{bridgebox}

\subsection{Explicit Hamiltonian Terms ($n=4$, $\lambda=1$)}

Applying equation~\eqref{eq:H-portfolio} to our $Q$ matrix yields
the following Pauli terms and constant (the same numbers the companion
notebook prints):

\begin{center}
\begin{tabular}{lrl}
    \toprule
    Pauli term & Coefficient & Origin \\
    \midrule
    $Z_0 Z_1$ & $+1.5294$ & $Q_{01}/2$ (AAPL--MSFT covariance + penalty) \\
    $Z_0 Z_2$ & $+1.5156$ & $Q_{02}/2$ (AAPL--JPM covariance + penalty) \\
    $Z_0 Z_3$ & $+1.5102$ & $Q_{03}/2$ (AAPL--XOM covariance + penalty) \\
    $Z_1 Z_2$ & $+1.5136$ & $Q_{12}/2$ (MSFT--JPM covariance + penalty) \\
    $Z_1 Z_3$ & $+1.5085$ & $Q_{13}/2$ (MSFT--XOM covariance + penalty) \\
    $Z_2 Z_3$ & $+1.5147$ & $Q_{23}/2$ (JPM--XOM covariance + penalty) \\
    $Z_0$     & $+0.0728$ & $-\tfrac12\sum_j Q_{0j}$ (AAPL) \\
    $Z_1$     & $+0.0471$ & $-\tfrac12\sum_j Q_{1j}$ (MSFT) \\
    $Z_2$     & $+0.0022$ & $-\tfrac12\sum_j Q_{2j}$ (JPM) \\
    $Z_3$     & $+0.0154$ & $-\tfrac12\sum_j Q_{3j}$ (XOM) \\
    \midrule
    Constant  & $-9.2294$ & dropped: $x^TQx = H_\text{portfolio} - 9.2294$ \\
    \bottomrule
\end{tabular}
\end{center}

\noindent Ground-state energy: $-3.1060 = -12.3354 + 9.2294$ (AAPL + MSFT).

\begin{defbox}[QAOA in one box]
\textbf{QAOA} --- Quantum Approximate Optimization Algorithm (developed in
full in Lab 4; everything needed here is below).
\begin{enumerate}
    \item \textbf{Start:} a Hadamard gate on every qubit gives all $2^4 = 16$
          bit strings the same probability, $1/16$.
    \item \textbf{Cost layer} $e^{-i\gamma H_\text{portfolio}}$: gives every bit
          string a phase proportional to its cost; $\gamma$ is an angle we
          choose. A phase alone changes no measurement probability. Each
          $Z_iZ_j$ term is built as CNOT--$R_z$--CNOT: the first CNOT writes
          ``do bits $i$ and $j$ differ?'' onto qubit $j$, the $R_z$ applies a
          phase that depends on that answer, and the second CNOT undoes the
          first. Each single $Z_i$ term is one $R_z$ on qubit $i$.
    \item \textbf{Mixer layer}: an $R_x(2\beta)$ rotation on every qubit makes
          bit strings that differ by one flip interfere, and this turns phase
          differences into probability differences.
    \item \textbf{Repeat} cost + mixer $p$ times (here $p = 1$ or $2$):
          $2p$ angles in total.
    \item \textbf{Tune:} run the circuit many times, average the cost of the
          outcomes (this average is written $\expect{H}$ and called the
          \emph{expectation value}), and let an ordinary classical optimizer
          (COBYLA, a standard derivative-free routine) change the angles to
          lower that average. Nothing in this step is quantum.
\end{enumerate}
The hope is that cheap bit strings end up with high probability. Whether
that happens here is what Part D reports.
\end{defbox}

\subsection{QAOA Circuit Structure ($p=1$)}

The $p=1$ QAOA circuit alternates one cost unitary and one mixer:

\begin{equation}
    \ket{\gamma,\beta} =
    e^{-i\beta\sum_i X_i}\,
    e^{-i\gamma H_\text{portfolio}}\,
    H^{\otimes 4}\ket{0000}.
\end{equation}

Per cost layer: one CNOT--$R_z$--CNOT block for each of the
$\binom{4}{2} = 6$ $Z_iZ_j$ terms (12 CNOTs, 6 $R_z$), plus one $R_z$ for
each of the 4 $Z_i$ terms; the mixer adds one $R_x$ per qubit. The
two-qubit part grows as $\binom{n}{2} = \mathcal{O}(n^2)$.

% ══════════════════════════════════════════════════════════════════════════════
\section{Part D: Result Interpretation}
% ══════════════════════════════════════════════════════════════════════════════

\subsection{Reading the Output Distribution}

After QAOA optimisation, we sample the circuit $M$ times and obtain a
histogram over $2^n = 16$ bitstrings. The \emph{technical} interpretation:
the cheapest portfolio, AAPL + MSFT, is the state $\ket{1100}$ (lab order
$\ket{q_0q_1q_2q_3}$, convention A1; Qiskit prints \texttt{0011}). Ideally it
would carry the highest probability; the next subsection shows that here it
does not.

\begin{finbox}[Technical $\to$ Business Translation]
\begin{center}
\begin{tabular}{cl}
    \toprule
    Bitstring & Business meaning \\
    \midrule
    $\ket{1100}$ & Invest in AAPL and MSFT \\
    $\ket{1001}$ & Invest in AAPL and XOM \\
    $\ket{0011}$ & Invest in JPM and XOM \\
    $\ket{0000}$ & Infeasible (no assets) \\
    $\ket{1111}$ & Infeasible (all assets, $K \neq 2$) \\
    \bottomrule
\end{tabular}
\end{center}
Infeasible states have high energy due to the penalty term: at $p=1$
they still carry about 26\% of the probability, at $p=2$ almost none.
\end{finbox}

\subsection{Approximation Quality and Business Risk}

Define the \emph{success probability} $P^* = P(\ket{1100})$: the
probability the circuit returns the exact optimal portfolio. With the angles
optimised in a noiseless simulation (single-$Z$ terms included):
\begin{itemize}
    \item $p=1$: about 74\% of the probability sits on the six feasible
          portfolios, each $\approx 12\%$, so $P^* \approx 12\%$; about 26\%
          is still infeasible.
    \item $p=2$: essentially 100\% is feasible, spread almost evenly
          ($\approx 16.7\%$ each, $P^* \approx 16.5\%$).
\end{itemize}
QAOA enforces the budget constraint, but it does not pick the winner: the six
feasible costs differ by at most $0.21$ and the best two by only $0.023$, far
below what one or two layers can resolve. The quantum step supplies a
\emph{shortlist}; the exact cost ranks it. For production use, we would:

\begin{enumerate}
    \item Sample the circuit $M$ times (e.g.\ $M=1000$).
    \item Keep the \emph{feasible} bitstrings that appear (the shortlist).
    \item Rank them by their exact cost $C(x)$, not by how often they were
          measured.
\end{enumerate}

Increasing $p$ raises the feasible fraction but also increases circuit
depth and susceptibility to hardware noise.

\subsection{Sensitivity Analysis: Risk Aversion $\lambda$}

\begin{mathbox}[How $\lambda$ Shifts the Optimal Portfolio]
At $\lambda = 0.5$ (return-focused): $\cblu{\text{AAPL}+\text{MSFT}}$ optimal. \\
At $\lambda = 1.0$ (balanced):      $\cblu{\text{AAPL}+\text{MSFT}}$ optimal. \\
At $\lambda = 2.0$ (risk-averse):   $\corg{\text{AAPL}+\text{JPM}}$ takes over. \\
At $\lambda = 5.0$ (very defensive): $\crd{\text{MSFT}+\text{JPM}}$ (lowest-variance
                                                                    pair) wins.

\noindent The quantum circuit must be re-optimised for each $\lambda$
because the Hamiltonian $H_\text{portfolio}(\lambda)$ changes.
The optimal QAOA parameters $(\gamma^*,\beta^*)$ also shift.
\end{mathbox}

% ══════════════════════════════════════════════════════════════════════════════
\section{Part E: IBM Composer Exercise}
% ══════════════════════════════════════════════════════════════════════════════

\textit{Platform: IBM Quantum Composer ---
\url{https://quantum.ibm.com/composer}}
% CHECK-FACT D4: platform URL and UI views not verified

\begin{studentbox}
\textbf{Build one full QAOA layer ($p=1$) by hand.}
The six $Z_iZ_j$ coefficients are equal to within 1.4\% ($1.5085$--$1.5294$,
see the Hamiltonian table in Part C), so one angle serves them all. The four
$Z_i$ terms are at most 5\% of a $ZZ$ term and are omitted here: this is a
\textbf{declared simplification}, and it changes the feasible fraction by less
than 0.5 percentage points.
\begin{enumerate}
    \item Put $H$ on $q_0, \dots, q_3$.
    \item For each of the 6 pairs $(q_i, q_j)$, $i < j$, add
          CNOT$(q_i, q_j)$ --- $R_z(0.50)$ on $q_j$ --- CNOT$(q_i, q_j)$.
          These 18 gates are the cost layer with $\gamma \approx 0.16$.
    \item Add $R_x(-0.56)$ on every qubit. This is the mixer with
          $\beta \approx -0.28$. If the Composer accepts only positive
          angles, use $5.72$.
    \item Open the probabilities view. The six bit strings with exactly two
          1s now carry $\approx 74\%$ in total ($\approx 12\%$ each).
\end{enumerate}
Compare with two variants:
\begin{itemize}
    \item \textbf{No cost layer} (only $H$, then $R_x$): $37.5\% = 6/16$.
          The mixer alone does nothing useful.
    \item \textbf{$R_x(+0.56)$} instead of $R_x(-0.56)$: $\approx 18\%$. The
          same gates with the opposite sign push probability \emph{away}
          from the feasible portfolios.
\end{itemize}
\textbf{Observation:} the angles are the whole algorithm, and that is why a
classical optimizer has to tune them.
\end{studentbox}

% ══════════════════════════════════════════════════════════════════════════════
\section{Part F: Qiskit Python Implementation}
% ══════════════════════════════════════════════════════════════════════════════

The full implementation is in the companion notebook. Three exercises
are provided:

\begin{enumerate}
    \item \textbf{QUBO Construction:} Build $Q$ from real asset data,
          enumerate all $\binom{4}{2}$ portfolios classically, identify
          the optimum.
    \item \textbf{QAOA Optimisation:} Run QAOA with $p \in \{1,2\}$,
          compare the measured distribution against brute-force.
    \item \textbf{Business Analysis:} Plot the efficient frontier, rank
          the feasible portfolios in the QAOA sample by exact cost, and
          compute the risk-return profiles of the top three.
\end{enumerate}

\begin{tcolorbox}[colback=green!5!white,colframe=green!75!black,
    title=Key Code Structure, fonttitle=\bfseries]
\begin{lstlisting}[language=Python]
# Core objects: same structure as Lab 4, different Hamiltonian
def build_portfolio_hamiltonian(Q: np.ndarray) -> SparsePauliOp:
    n     = Q.shape[0]
    terms = []
    for i in range(n):
        for j in range(i+1, n):
            # zz_str, z_str: Pauli-label helpers, see notebook
            coeff = Q[i,j] / 2           # ZiZj coefficient
            terms.append((zz_str(i,j,n), coeff))
        z_coeff = -sum(Q[i,:]) / 2       # Zi coefficient
        terms.append((z_str(i,n), z_coeff))
    return SparsePauliOp.from_list(terms).simplify()

# QAOA circuit: H layer, then per layer one CNOT-Rz-CNOT per ZZ term,
# one Rz per Z term, Rx mixer (see notebook)
\end{lstlisting}
\end{tcolorbox}

% ══════════════════════════════════════════════════════════════════════════════
\section{Part G: Scaling and Outlook}
% ══════════════════════════════════════════════════════════════════════════════

\subsection{Scaling to Larger Portfolios}

\begin{table}[H]
\centering
\caption{Portfolio counts and QAOA circuit size ($n$ assets, $K$ selected;
one qubit per asset).}
\label{tab:scaling}
\begin{tabular}{ccccc}
    \toprule
    $n$ & $K$ & Portfolios $\binom{n}{K}$ & ZZ terms $\binom{n}{2}$ & CNOTs per cost layer \\
    \midrule
    4  & 2  & 6                         & 6     & 12    \\
    6  & 3  & 20                        & 15    & 30    \\
    10 & 5  & 252                       & 45    & 90    \\
    20 & 10 & 184{,}756                 & 190   & 380   \\
    50 & 10 & $\approx 1.0 \times 10^{10}$ & 1{,}225 & 2{,}450 \\
    50 & 25 & $\approx 1.3 \times 10^{14}$ & 1{,}225 & 2{,}450 \\
    \bottomrule
\end{tabular}
\end{table}

The QAOA circuit grows as $\mathcal{O}(n^2)$ gates (from the $\binom{n}{2}$
ZZ interaction terms), while the classical search space grows as
$\binom{n}{K} \sim n^K / K!$. For $n=20$, $K=5$: classical brute-force
checks $\binom{20}{5} = 15{,}504 \approx 1.6 \times 10^4$ portfolios (fast
classically); for $n=50$, $K=10$: $\approx 1.0 \times 10^{10}$ (impractical
to enumerate). The QAOA circuit for $n=50$ needs $\binom{50}{2} = 1225$
CNOT--$R_z$--CNOT blocks, i.e.\ $2450$ CNOTs per layer: small in gate count,
but see the noise estimate below.

\subsection{Limitations and Open Questions}

\begin{itemize}
    \item \textbf{Barren plateaus:} For $n \gtrsim 50$, the QAOA landscape
          may be exponentially flat, requiring problem-informed initial
          parameters.
    \item \textbf{Hardware noise:} The $2450$ CX gates of one layer at
          $n=50$, with $\sim\!0.5\%$ error each, give an overall fidelity of
          $0.995^{2450} \approx 5 \times 10^{-6}$ — noise dominates.
    \item \textbf{Quantum advantage:} Not yet demonstrated for portfolio
          optimization. Active research area (JPMorgan, Goldman Sachs,
          IBM Research).
    \item \textbf{Beyond QUBO:} Transaction costs, lot sizes, turnover
          constraints are not naturally quadratic — require slack variables
          or penalty engineering.
\end{itemize}

% ══════════════════════════════════════════════════════════════════════════════
\section{Part H: Assessment}
% ══════════════════════════════════════════════════════════════════════════════

\textbf{Q1.\ QUBO Formulation}\\
Why is the budget constraint $\sum_i x_i = K$ added as a penalty term
$\lambda_\text{pen}(\sum_i x_i - K)^2$ rather than enforced directly in
the quantum circuit?
\begin{itemize}
    \item[A)] Because quantum circuits cannot measure the number of 1s
              in a bitstring.
    \item[B)] Because the QAOA Hamiltonian must be a sum of Pauli operators
              (diagonal in the $Z$-basis), and a soft penalty converts the
              hard constraint into such a form.
    \item[C)] Because the penalty term is always zero for feasible solutions,
              making it costless.
\end{itemize}

\vspace{0.3cm}
\textbf{Q2.\ Hamiltonian Terms}\\
In the portfolio Hamiltonian $H_\text{portfolio}$, what is the physical
origin of the $Z_iZ_j$ terms?
\begin{itemize}
    \item[A)] The expected returns $\mu_i$ of the individual assets.
    \item[B)] The covariance $\Sigma_{ij}$ between assets $i$ and $j$,
              plus the budget penalty terms that couple all pairs of qubits.
    \item[C)] The CNOT gates in the QAOA mixer.
\end{itemize}

\vspace{0.3cm}
\textbf{Q3.\ Result Interpretation}\\
At $p = 2$ the QAOA circuit for the 4-asset problem returns each of the
six feasible bit strings with probability $\approx 16$--$17\%$ and almost
nothing else. What should the analyst recommend?
\begin{itemize}
    \item[A)] The most frequent bit string in the sample.
    \item[B)] Among the feasible bit strings, the one with the lowest exact
              cost $C(x)$ --- AAPL + MSFT --- with the next-lowest
              (AAPL + XOM) as the alternative.
    \item[C)] Re-run with more shots until one bit string exceeds 99\%.
\end{itemize}

\vspace{0.3cm}
\textbf{Q4.\ Scaling}\\
For $n=50$ assets and $K=10$, the classical brute-force requires
$\binom{50}{10} \approx 10^{10}$ cost evaluations. How many CNOT gates
does one QAOA layer need for $n = 50$?
\begin{itemize}
    \item[A)] $n = 50$
    \item[B)] $n^2 = 2500$
    \item[C)] $2\binom{n}{2} = 2450$ (one CNOT pair per $Z_iZ_j$ term)
\end{itemize}

\vspace{1cm}
\hrule
\vspace{0.2cm}
\footnotesize{\textbf{Solutions:}
Q1: B (the Ising Hamiltonian is $\sum_i h_i Z_i + \sum_{ij} J_{ij} Z_iZ_j$;
    the penalty converts the constraint into this Pauli-operator form);
Q2: B (the $Q_{ij}$ for $i\neq j$ contains $\lambda\Sigma_{ij}$ from
    the risk term plus $\lambda_\text{pen}$ from the penalty, both
    contributing to the $Z_iZ_j$ coefficient after the substitution
    $x_i=(1-Z_i)/2$);
Q3: B (shot counts differ here only by sampling noise; QAOA supplied the
    feasible shortlist, and the exact cost ranks it);
Q4: C ($\binom{50}{2} = 1225$ $Z_iZ_j$ terms, each built as
    CNOT--$R_z$--CNOT, give $2\times 1225 = 2450$ CNOTs per layer, plus
    $1225 + 50$ single-qubit $R_z$ and $50$ $R_x$; total $\mathcal{O}(n^2)$).}

\end{document}